\documentclass{article}
\usepackage[T1]{fontenc}
\usepackage{lmodern}
\usepackage{graphicx}
\usepackage{amsmath,amssymb}
\usepackage[a4paper,margin=2cm]{geometry}
\usepackage[absolute,overlay]{textpos}
\setlength{\TPHorizModule}{1mm}
\setlength{\TPVertModule}{1mm}
\usepackage[indonesian]{babel}
\usepackage{hyperref}
\setlength{\emergencystretch}{2em}
% unit: Halaman 7

\begin{document}

% === Halaman 7 (python/native) ===

• Misalkan e dan d dipilih sedemikian sehingga 

  e  d  1 (mod (n))

 atau


e  d = k(n) + 1

 Maka


mk(n) + 1  m (mod n)






me d  m (mod n)  → 

(me)d  m (mod n)

• Enkripsi:  c = Ee(m) = me mod n

• Dekripsi:  m = Dd(c) = cd mod n 

7

\end{document}
